% Sesi 2: grafik akademik dan data fiktif.
\documentclass{article}
\usepackage[margin=2cm]{geometry}
\usepackage{pgfplots}
\pgfplotsset{compat=1.18}
\begin{document}
\begin{figure}[ht]\centering
\begin{tikzpicture}
\begin{axis}[width=11cm,height=7cm,xlabel={Waktu (minggu)},ylabel={Nilai (satuan)},
 xmin=0,xmax=5,ymin=0,ymax=24,grid=major,grid style={gray!20},legend pos=north west]
\addplot[blue,thick,mark=*] table[x=x,y=y,col sep=comma] {data.csv};
\addlegendentry{Pengamatan fiktif}
\addplot[black,dashed,domain=0:5,samples=2] {10+2*x};
\addlegendentry{Acuan}
\node[anchor=south east,align=right] at (axis cs:5,21) {Titik terakhir};
\end{axis}\end{tikzpicture}
\caption{Data contoh fiktif dibandingkan dengan kurva acuan.}
\end{figure}\end{document}
